This thesis investigates whether Reddit post popularity can be predicted using publicly available data and off-the-shelf machine learning tools, and whether such predictions can inform automated content generation and deployment. The work addresses both the technical feasibility of constructing such a pipeline and its adversarial implications for platform integrity.

The study frames popularity prediction as a regression task targeting upvote ratio -- a normalized measure of community approval less susceptible to exposure bias than raw engagement metrics. Using historical Reddit data from 2008--2022 across 19 politically diverse subreddits, XGBoost regression models were trained on Transformer-based text embeddings (Stella 400M) combined with structured metadata features including temporal patterns, author history, domain characteristics, and sentiment indicators. Offline cross-validation demonstrated moderate predictive performance, achieving Spearman correlations exceeding 0.20 and mean absolute errors around 7\% when predicting upvote ratios on held-out data.

To assess external validity, a live deployment experiment was conducted in which posts generated by a large language model were scored by the trained predictor and selectively published to Reddit. This experiment yielded a significant negative result: model-approved posts (top 10\% of predicted upvote ratios) substantially underperformed randomly selected posts, achieving mean upvote ratios approximately 75\% lower than the random baseline. Within the model-approved subset, predictions exhibited a strong negative correlation with actual outcomes ($\rho = -0.596$, $p < 0.0001$), indicating systematic rank inversion rather than weak signal. Conversely, randomly selected generated posts achieved mean upvote ratios exceeding historical training baselines, demonstrating that content generation succeeded while model-guided selection failed.

These findings reveal a critical distribution shift between offline validation and real-world deployment: models trained on human-authored posts embedded in organic community dynamics fail when applied to generated content from unfamiliar accounts. The deployment also encountered substantial moderation resistance, with account bans in four of nine targeted subreddits and individual post removals across remaining communities, highlighting social and procedural barriers beyond content quality.

The thesis demonstrates that while the technical barrier to assembling such a pipeline is low -- requiring only scraped data, public embedding models, and standard ML libraries -- offline predictive accuracy does not guarantee deployment effectiveness. The work contributes a popularity regression framework, systematic analysis of cross-subreddit generalization, and empirical documentation of deployment failure modes, with implications for understanding both the capabilities and limitations of automated influence operations on social media platforms.